% !TEX root = ../vista.tex
% EJEMPLO: borra este fichero cuando crees los tuyos (fca.tex, dl.tex...).

\begin{definicion}[Operador de cierre]{op-cierre}
    Un \emph{operador de cierre} sobre un conjunto $X$ es una aplicación
    $c\colon \Pot(X)\to\Pot(X)$ que es extensiva ($A\subseteq c(A)$),
    monótona ($A\subseteq B \Rightarrow c(A)\subseteq c(B)$) e
    idempotente ($c(c(A))=c(A)$).
\end{definicion}

\begin{lema}[Cierre de una unión]{cierre-union}
    Si $c$ es un operador de cierre, entonces
    $c(A\cup B)=c\bigl(c(A)\cup c(B)\bigr)$.
\end{lema}
\begin{proof}
    ($\subseteq$) Como $A\cup B\subseteq c(A)\cup c(B)$ por extensividad,
    la monotonía da la inclusión.
    ($\supseteq$) Por monotonía $c(A)\cup c(B)\subseteq c(A\cup B)$, y aplicando
    de nuevo monotonía e idempotencia,
    $c\bigl(c(A)\cup c(B)\bigr)\subseteq c(c(A\cup B))=c(A\cup B)$.
\end{proof}

\begin{idea}{idea-generalizar}
    Ver si el lema anterior se generaliza a uniones arbitrariasdfgfdg de familias.
    \pendiente{comprobar el caso de la familia infinita}
\end{idea}
